\section{Objective and Impersonal Knowledge}

Yesterday, we saw that \textbf{Julius Evola} claimed both knowledge and action as the two basic elements. Let's be clear about what is meant by knowledge. This is not the scholarly knowledge of the professors nor is it related to any discipline like comparative religion. Rather, this knowledge is a gnosis, that is, a direct, intuitive, more than personal knowing which brings absolute certainty. In the snippet translated below the horizontal line, taken from Chapter 2, Section 2, of Sintesi di dottrina della razza, Evola begins his explication of race as based on the traditional understanding of man. It is insufficient to know that the Greeks and Scholastics taught the same view. The one thing necessary is to know oneself as a being consisting of Spirit, Soul, and Body. A man must know and be able to disntinguish these parts of himself within his own consciousness. I know first hand that developing this knowledge is a difficult and arduous process, more difficult than any university class or perfecting a skill in combat or the arts. Even when this knowledge is attained, it is yet more difficult to maintain one's awareness as spirit on a consistent or permanent basis. Many speak of spirit; few are spirit. The risk of taking this path is the overturning of one's cherished beliefs. The risk of not taking it is to continue to live like an animal.

\paragraph{The Three Levels of the Doctrine of Race}

Since we want to clarify the doctrine of race from the traditional point of view, of course we will assume, as a premise, the traditional conception of the human being, according to which, man as such cannot be reduced to purely biological, instinctive, hereditary, or naturalistic determinisms: if all that has one of its parts, ignored by a dubious spiritualism, exaggerated by a myopic positivism, it still remains a fact that man is distinguished from animals insofar as he participates in a supernatural, super-biological element, by virtue of which he can be free or himself. Between the one and the other, as an intermediate element, in a certain way he remains an animal. The distinction in the human being of three different principles, namely, the body, soul, and spirit, is fundamental for the traditional view. In a more or less complete form, it is found in the teachings of all the ancient traditions, and it continued up to the Middle Ages:

\begin{itemize}


\item The Aristotelian and Scholastic conception of the three souls: vegetative, sensitive, and intellectual

\item In the Hellenic trinity of \emph{soma}, \emph{psyche}, and \emph{nous}

\item The roman conception of \emph{mens}, \emph{anima}, and \emph{corpus}

\item The Indo-Aryan of \emph{sthula-sarira}, \emph{linga-sarira} and \emph{karana-sarira}

\end{itemize}


And so on: these are all equivalent expressions of it. It is important to emphasize that this view must not be considered as one particular philosophical interpretation among so many others, to be discussed, criticized, or debated, but as an objective and impersonal knowledge, conforming to the very nature of things.


\flrightit{Posted on 2011-08-11 by Aeneas}

